\section*{3 The Top-Down Approach}
Our two main reductions both take a ``top-down'' approach to finding an optimal formula. That is, given a function $f$, they try to find functions $g$ and $h$ such that $g$ and $h$ are the two functions fed into the final output gate in an optimal formula for $f$ and then recursing.
This is formalized as follows.
\setcounter{nttheorem}{16}
\begin{ntdefinition}[Author Definition 17: Optimal Subcomputations Set]
Let $f:\{0,1\}^n\to\{0,1\}$. We define the set of optimal subcomputations for $f$, denoted $\Opt(f)$, as follows.
Let $g,h:\{0,1\}^n\to\{0,1\}$ be Boolean functions of the same arity as $f$ and $O\in\{\land,\lor\}$. Then $\{g,O,h\}\in\Opt(f)$ if and only if there exists an optimal formula $\phi=\phi_gO\phi_h$ for computing $f$ such that $\phi_g$ computes $g$ and $\phi_h$ computes $h$.
\end{ntdefinition}
We note that in this definition we are implicitly using that the gate set $\{\land,\lor\}$ is symmetric with respect to its inputs.
We say a function $g$ is in an optimal subcomputation for $f$ if $g$ is contained in some element of $\Opt(f)$. In other words, $g$ is in an optimal subcomputation for $f$ if there exists an $h$ and $O$ such that $\{g,O,h\}\in\Opt(f)$.
It is easy to see that $\Opt(f)$ is almost always non-empty.
\begin{ntproposition}[Author Proposition 18]
Let $f:\{0,1\}^n\to\{0,1\}$ such that $L(f)\geq 2$. Then $\Opt(f)$ is non-empty.
\end{ntproposition}
Next, we can define the problem of finding an optimal subcomputation.
\begin{ntdefinition}[Author Definition 19: Decomposition Problem]
The Decomposition Problem, DecompProblem is as follows:
\begin{description}\item[Given:] the truth table of a Boolean function $f$ satisfying $L(f)\geq 2$\item[Output:] some element of $\Opt(f)$.\end{description}
\end{ntdefinition}
It is easy to see that DecompProblem is equivalent to Search-MFSP. DecompProblem can be easily solved with an oracle to Search-MFSP. The following recursive procedure shows the reverse direction.
\begin{nttheorem}[Author Theorem 20: Search-MFSP reduces to DecompProblem]
There is a deterministic $O(N^2)$-time algorithm for solving Search-MFSP on inputs of length $N$ given access to an oracle that solve DecompProblem on instances of length $N$.
\end{nttheorem}
\begin{proof}
The pseudocode for this reduction is written in Algorithm 1.
\begin{quote}\textbf{Algorithm 1. Reduction from Search-MFSP to DecompProblem.}\par
\textbf{procedure FindOptFormula($f$)}\par
Given the length-$N$ truth table of a function $f$ that takes $n$-inputs and oracle access to DecompProblem return an optimal formula for $f$.
\begin{enumerate}
\item If there exists a size one formula $\phi$ computing $f$, then return $\phi$. End if.
\item Let $\{g,O,h\}$ be the output returned by the oracle $\operatorname{DecompProblem}(f)$.
\item Recursively compute the formula $\phi_g\leftarrow\operatorname{FindOptFormula}(g)$.
\item Recursively compute the formula $\phi_h\leftarrow\operatorname{FindOptFormula}(h)$.
\item Return the formula given by $\phi_gO\phi_h$.
\end{enumerate}
\textbf{end procedure}
\end{quote}
The correctness of this algorithm is easy to see as long as one is able to bound the number of recursive calls the algorithm makes. To see that the number of recursive calls is bounded by $O(N)$, notice that each iteration of the algorithm reveals one more gate in the optimal formula for $f$. Thus, since $L(f)=O(N)$, we have that there are at most $O(N)$ recursive calls.
\end{proof}
Our goal is now to try to solve DecompProblem (i.e. find an element of $\Opt(f)$) given an oracle to MFSP. Recall from the introduction that our high-level approach is as follows
\begin{enumerate}
\item Find an efficient ``test'' that functions that in an optimal subcomputation of $f$ pass but not too many other functions pass.
\item Efficiently build the (not too long) list Candidates of things that pass the test.
\item Iterate through all pairs of elements in Candidates and all possible gates, and efficiently check if this yields an element of $\Opt(f)$.
\end{enumerate}
Item 1 will be the subject of Section 4, Item 2 will be different in our two main reductions, and Item 3 is provided by the next lemma.
\begin{ntlemma}[Author Lemma 21: Test membership in $\Opt(f)$ efficiently with MFSP]
Let $f,g,h:\{0,1\}^n\to\{0,1\}$. Then
\[\{g,O,h\}\in\Opt(f)\quad\Longleftrightarrow\quad f=gOh\text{ and }L(f)=L(g)+L(h).\]
\end{ntlemma}
\begin{proof}
We prove the forward direction first. Suppose that $\{g,O,h\}\in\Opt(f)$. Then there exists an optimal formula $\phi=\phi_gO\phi_h$ for computing $f$ such that $\phi_g$ computes $g$ and $\phi_h$ computes $h$. Clearly this implies that $f=gOh$.
Moreover, $|\phi|=|\phi_g|+|\phi_h|$. On the other hand, since $\phi$ is optimal, we have that $|\phi|=L(f)$, $|\phi_g|=L(g)$, and $|\phi_h|=L(h)$. (Otherwise, one could build a smaller formula for $f$ by replacing $\phi_g$ or $\phi_h$ with a smaller formula computing the same function.) Hence $L(f)=L(g)+L(h)$.
For the reverse direction, suppose that $L(f)=L(g)+L(h)$ and $f=gOh$. Let $\phi_g$ and $\phi_h$ be optimal formulas for $g$ and $h$. Then $\phi=\phi_gO\phi_h$ clearly computes $f$ and has size $L(f)$. Hence $\{g,O,h\}\in\Opt(f)$.
\end{proof}
